%%%%%%%%%%%%%%%%%%%%%%%%%%%%%%%%%%%%%%%%%%  不使用 authblk 包制作标题  %%%%%%%%%%%%%%%%%%%%%%%%%%%%%%%%%%%%%%%%%%%%%%
%-------------------------------PPT Title-------------------------------------
\title{\textrm{AI for Science:~}数据驱动下的\材料模拟与技术服务}
%\subtitle{两弹元勋与核弹原理}
%-----------------------------------------------------------------------------

%----------------------------Author & Date------------------------------------
%\author[]{\vskip +10pt 姜\;\;骏\inst{}} %[]{} (optional, use only with lots of authors)
%% - Give the names in the same order as the appear in the paper.
%% - Use the \inst{?} command only if the authors have different
%%   affiliation.
\institute[BCC]{\inst{}%
%\institute[Gain~Strong]{\inst{}%
\vskip -5pt 北京市计算中心有限公司~材料计算团队}
%\vskip -20pt {\large 格致斯创~科技}}
\date[\today] % (optional, should be abbreviation of conference name)
{	%{\fontsize{6.2pt}{4.2pt}\selectfont{\textcolor{blue}{E-mail:~}\url{jiangjun@bcc.ac.cn}}}
\vskip 25 pt {\fontsize{8.2pt}{6.2pt}\selectfont{%清华大学\;\;物理系% 报告地点
	\vskip 5 pt \textrm{2026.~04.~17}}}
}

%% - Either use conference name or its abbreviation
%% - Not really information to the audience, more for people (including
%%   yourself) who are reading the slides onlin%%   yourself) who are reading the slides onlin%%   yourself) who are reading the slides onlineee
%%%%%%%%%%%%%%%%%%%%%%%%%%%%%%%%%%%%%%%%%%%%%%%%%%%%%%%%%%%%%%%%%%%%%%%%%%%%%%%%%%%%%%%%%%%%%%%%%%%%%%%%%%%%%%%%%%%%%

\subject{}
% This is only inserted into the PDF information catalog. Can be left
% out.
%\maketitle
\frame
{
%	\frametitle{\fontsize{9.5pt}{5.2pt}\selectfont{\textcolor{orange}{“高通量并发式材料计算算法与软件”年度检查}}}
\titlepage
}
%-----------------------------------------------------------------------------

%------------------------------------------------------------------------------列出全文 outline ---------------------------------------------------------------------------------
\section*{}
\frame[allowframebreaks]
{
	\frametitle{\textrm{Outline}}
%  \frametitle{\textcolor{mycolor}{\secname}}
  \tableofcontents%[current,currentsection,currentsubsection]
}
%在每个section之前列出全部Outline
%类似的在每个subsection之前列出全部Outline是\AtBeginSubsection[]
%\AtBeginSection[]
%{
%  \frame<handout:0>%[allowframebreaks]
%  {
%    \frametitle{Outline}
%%全部Outline中，本部分加亮
%    \tableofcontents[current,currentsection]
%  }
%}

%-----------------------------------------------PPT main Body------------------------------------------------------------------------------------
\small
\section{引言与平台概览}

\begin{frame}{课程目标}
\begin{itemize}
    \item 了解当前主流智能体开发平台
    \item 掌握扣子(Coze)平台的完整功能
    \item 学会使用n8n构建自动化工作流
    \item 对比不同平台的适用场景
\end{itemize}
\end{frame}

\begin{frame}{智能体平台生态}
\begin{center}
\framebox{\parbox{0.8\textwidth}{低代码/无代码智能体开发平台正在爆发}}
\end{center}
\begin{itemize}
    \item 降低AI应用开发门槛
    \item 提供可视化编排、预置节点、模板库
    \item 支持RAG、函数调用、多智能体协作
\end{itemize}
\end{frame}

\begin{frame}{常用平台工具一览}
\begin{table}
\small
\begin{tabular}{lll}
\toprule
平台 & 类型 & 特点 \\
\midrule
扣子(Coze) & 国内主流 & 字节出品，功能全面 \\
trae & 代码辅助 & 智能体驱动IDE \\
Clawdbot & 自动化 & 爬虫与数据处理 \\
Skill & 技能平台 & 预置技能市场 \\
Manus & 通用智能体 & 任务执行 \\
n8n & 工作流 & 开源自动化 \\
ChatGPT Agent & 官方 & OpenAI智能体 \\
cursor & 代码编辑器 & AI原生IDE \\
Dify & 开源LLMOps & 工作流+RAG \\
Dokie & 文档智能 & 知识库问答 \\
\bottomrule
\end{tabular}
\end{table}
\end{frame}

% 每个平台简要介绍（2页/个，共10个平台，20页，但节省篇幅，每个1页）
\begin{frame}{扣子(Coze) – 字节跳动智能体平台}
\begin{itemize}
    \item 免费、中文友好、插件丰富
    \item 支持知识库、工作流、多模型
    \item 可发布到飞书、微信、Discord
\end{itemize}
\end{frame}

\begin{frame}{trae – 智能体驱动IDE}
\begin{itemize}
    \item 类似Cursor，内置AI编程智能体
    \item 可自动生成代码、修复Bug
    \item 与项目上下文深度集成
\end{itemize}
\end{frame}

\begin{frame}{Clawdbot – 爬虫与数据自动化}
\begin{itemize}
    \item 智能体负责数据采集、清洗
    \item 支持动态网站、反爬策略
    \item 可与n8n联动
\end{itemize}
\end{frame}

\begin{frame}{Skill – 智能体技能平台}
\begin{itemize}
    \item 预置技能市场（天气、新闻、邮件）
    \item 支持自定义技能上传
    \item 跨智能体共享
\end{itemize}
\end{frame}

\begin{frame}{Manus – 通用任务智能体}
\begin{itemize}
    \item 输入自然语言目标，自动规划执行
    \item 调用浏览器、代码、API
    \item 类似AutoGPT但更稳定
\end{itemize}
\end{frame}

\begin{frame}{n8n – 开源自动化工作流}
\begin{itemize}
    \item 节点化编排，支持400+应用
    \item 可自托管，数据隐私可控
    \item 内置AI节点（LangChain、OpenAI）
\end{itemize}
\end{frame}

\begin{frame}{ChatGPT Agent – OpenAI官方}
\begin{itemize}
    \item 通过Function Calling实现工具调用
    \item 支持代码解释器、文件检索
    \item 可定制指令
\end{itemize}
\end{frame}

\begin{frame}{cursor – AI原生代码编辑器}
\begin{itemize}
    \item 基于VSCode，深度集成GPT-4
    \item 智能体辅助编程、重构、测试
    \item 支持终端命令执行
\end{itemize}
\end{frame}

\begin{frame}{Dify – 开源LLM应用平台}
\begin{itemize}
    \item 可视化RAG流水线
    \item 支持工作流编排、API发布
    \item 私有化部署友好
\end{itemize}
\end{frame}

\begin{frame}{Dokie – 文档智能问答}
\begin{itemize}
    \item 上传文档，构建知识库智能体
    \item 支持PDF、Word、网页
    \item 可嵌入网页或Slack
\end{itemize}
\end{frame}

\section{扣子(Coze)平台深度解析}

\begin{frame}{扣子平台架构}
\begin{center}
\framebox{\parbox{0.7\textwidth}{智能体编辑界面（左侧插件、中间提示词、右侧预览）}}
\end{center}
\begin{itemize}
    \item 核心组件：提示词、知识库、插件、工作流、记忆
    \item 发布渠道：API、飞书、微信、Discord
\end{itemize}
\end{frame}

% 智能体创建
\begin{frame}{智能体创建 – 第一步}
\begin{enumerate}
    \item 登录扣子官网 (coze.cn)
    \item 点击“创建智能体”
    \item 输入名称、描述、头像
    \item 选择基础模型（云雀、智谱、月之暗面等）
\end{enumerate}
\end{frame}
\begin{frame}{智能体创建 – 人设与回复逻辑}
\begin{block}{人设提示词示例}
“你是一位资深金融分析师，回答风格专业、简洁，使用数据支撑观点。”
\end{block}
\begin{itemize}
    \item 设定角色、语气、禁止事项
    \item 可配置开场白、预置问题
\end{itemize}
\end{frame}

% 提示词优化
\begin{frame}{提示词优化 – 技巧}
\begin{itemize}
    \item 使用变量 `{{input}}` 接收用户输入
    \item 添加约束（字数、格式、引用来源）
    \item 提供少样本示例（Few-shot）
    \item 测试不同模型效果
\end{itemize}
\end{frame}
\begin{frame}{提示词优化 – 调试面板}
\begin{itemize}
    \item 实时预览回复
    \item 查看模型调用日志
    \item 调整温度、top\_p等参数
    \item A/B测试不同提示词版本
\end{itemize}
\end{frame}

% 工具调用
\begin{frame}{工具调用 – 插件市场}
\begin{itemize}
    \item 内置插件：搜索引擎、图片理解、天气、新闻
    \item 第三方插件：飞书文档、高德地图、GitHub
    \item 自定义插件：通过OpenAPI规范导入
\end{itemize}
\end{frame}
\begin{frame}{工具调用 – 工作流集成}
\begin{itemize}
    \item 智能体可调用已搭建的工作流
    \item 工作流节点包括：代码、HTTP请求、LLM、条件判断
    \item 实现复杂多步骤任务
\end{itemize}
\end{frame}

% 知识库构建
\begin{frame}{知识库构建 – 数据导入}
\begin{itemize}
    \item 支持格式：TXT、PDF、Word、Excel、网页链接
    \item 自动分块、向量化
    \item 支持手动分段、标签管理
\end{itemize}
\end{frame}
\begin{frame}{知识库构建 – 检索增强}
\begin{itemize}
    \item 智能体自动从知识库检索相关内容
    \item 可设置检索数量、相似度阈值
    \item 支持混合检索（关键词+向量）
\end{itemize}
\end{frame}

% 工作流搭建
\begin{frame}{工作流搭建 – 可视化编排}
\begin{center}
\framebox{\parbox{0.6\textwidth}{拖拽节点，连线编排}}
\end{center}
\begin{itemize}
    \item 开始节点：接收输入参数
    \item 大模型节点：调用LLM处理
    \item 代码节点：运行Python/JS
    \item 结束节点：输出结果
\end{itemize}
\end{frame}
\begin{frame}{工作流搭建 – 变量与数据流}
\begin{itemize}
    \item 支持全局变量、节点间传递
    \item 可使用 `{{node.output}}` 引用
    \item 条件分支、循环节点
\end{itemize}
\end{frame}

% 智能体调试
\begin{frame}{智能体调试 – 预览与日志}
\begin{itemize}
    \item 右侧对话框实时测试
    \item 查看每次工具调用详情
    \item 定位知识库检索结果
    \item 导出对话记录
\end{itemize}
\end{frame}
\begin{frame}{智能体调试 – 版本管理}
\begin{itemize}
    \item 保存为草稿或发布版本
    \item 回滚到历史版本
    \item 多环境（开发、测试、生产）
\end{itemize}
\end{frame}

\subsection{扣子平台工具节点详解}

\begin{frame}{搜索引擎工具节点}
\begin{itemize}
    \item 调用必应/百度搜索API
    \item 返回网页标题、摘要、链接
    \item 可设置搜索数量、语言、时间范围
\end{itemize}
\end{frame}
\begin{frame}{搜索引擎节点 – 使用示例}
```json
{
  "query": "2025年AI发展趋势",
  "count": 5,
  "language": "zh-CN"
}
```

输出：搜索结果列表，供LLM总结。
\end{frame}

\begin{frame}{文献库工具节点}
\begin{itemize}
\item 检索学术论文（arXiv、知网等）
\item 支持按标题、作者、年份筛选
\item 返回论文摘要、PDF链接
\end{itemize}
\end{frame}
\begin{frame}{文献库节点 – 应用场景}
智能体可充当科研助手：根据用户问题检索相关论文，并生成文献综述。
\end{frame}

\begin{frame}{图片理解工具节点}
\begin{itemize}
\item 调用多模态模型（如GPT-4V、Claude 3）
\item 输入图片URL或Base64
\item 输出图片描述、物体检测、OCR结果
\end{itemize}
\end{frame}
\begin{frame}{图片理解节点 – 示例}
用户上传产品图片，智能体识别品牌、型号、颜色，并推荐购买链接。
\end{frame}

\begin{frame}{图像处理程序节点}
\begin{itemize}
\item 无需编程，预置图像处理功能
\item 裁剪、缩放、滤镜、水印
\item 格式转换（PNG ↔ JPEG）
\item 人脸模糊、背景替换
\end{itemize}
\end{frame}
\begin{frame}{图像处理节点 – 工作流集成}
可与其他节点串联：图片理解 → 图像处理 → 存储到云盘 → 发送邮件。
\end{frame}

\begin{frame}{多模态大模型节点}
\begin{itemize}
\item 支持GPT-4V、Gemini Pro Vision等
\item 输入图文混合消息
\item 输出文本、JSON结构化数据
\item 适用于视觉问答、图表分析
\end{itemize}
\end{frame}
\begin{frame}{多模态节点 – 案例}
上传销售报表截图，提问“哪个季度增长率最高？”，模型直接给出答案。
\end{frame}

% 扣子部分结束，总页数约60页，现在进入n8n

\section{n8n平台功能介绍}

\begin{frame}{n8n 概览}
\begin{itemize}
\item 开源、自托管的自动化工作流引擎
\item 节点式编排，支持400+应用（Google Sheets, Slack, MySQL, HTTP等）
\item 内置AI节点（OpenAI、LangChain、Hugging Face）
\item 可定时触发、Webhook触发
\end{itemize}
\end{frame}

\begin{frame}{n8n 安装方式}
\begin{itemize}
\item Docker一键安装：docker run -it --rm --name n8n -p 5678:5678 n8nio/n8n
\item npm安装：npm install n8n -g
\item 云服务：n8n.cloud
\end{itemize}
\end{frame}

% 智能体创建
\begin{frame}{n8n中的“智能体”概念}
n8n本身是工作流引擎，但可结合AI节点构建“智能体”：
\begin{itemize}
\item 使用OpenAI节点实现意图识别
\item 条件分支动态选择下一步动作
\item 循环节点实现多步推理
\end{itemize}
\end{frame}
\begin{frame}{n8n智能体创建示例}
工作流：Webhook接收消息 → OpenAI分类意图 → 根据意图调用不同API → 返回结果。
\end{frame}

% 工作流创建
\begin{frame}{n8n工作流创建 – 界面}
\begin{center}
\framebox{\parbox{0.7\textwidth}{左侧节点库，中间画布，右侧节点配置}}
\end{center}
\begin{itemize}
\item 拖拽节点至画布
\item 连接节点（连线顺序）
\item 测试运行、查看执行日志
\end{itemize}
\end{frame}
\begin{frame}{工作流创建 – 触发方式}
\begin{itemize}
\item 定时触发（Cron）
\item Webhook触发（HTTP请求）
\item 手动触发（按钮）
\item 其他节点触发（如邮件接收）
\end{itemize}
\end{frame}

% 数据分析节点
\begin{frame}{数据分析节点 – 聚合}
\begin{itemize}
\item 对输入数据分组、计算总和/平均值
\item 支持SQL风格的聚合函数
\end{itemize}
\end{frame}
\begin{frame}{数据分析节点 – 排序与过滤}
\begin{itemize}
\item 按字段排序
\item 条件过滤（大于、包含、正则）
\item 分页限制
\end{itemize}
\end{frame}

% 大模型节点
\begin{frame}{大模型节点 – OpenAI}
\begin{itemize}
\item 配置API Key、模型（gpt-4o, gpt-3.5-turbo）
\item 输入消息列表（system, user, assistant）
\item 支持函数调用（Function Calling）
\end{itemize}
\end{frame}
\begin{frame}{大模型节点 – LangChain}
\begin{itemize}
\item 内置多种链（LLMChain, RetrievalQA）
\item 可连接向量数据库
\item 支持自定义提示词模板
\end{itemize}
\end{frame}
\begin{frame}{大模型节点 – 本地模型}
\begin{itemize}
\item 通过HTTP Request节点调用Ollama、LocalAI
\item 支持Llama、Mistral等
\end{itemize}
\end{frame}

% 接口调用节点
\begin{frame}{HTTP Request节点}
\begin{itemize}
\item 支持GET、POST、PUT、DELETE
\item 自定义Headers、Body（JSON/Form）
\item 支持OAuth2认证
\item 可解析响应JSON/XML
\end{itemize}
\end{frame}
\begin{frame}{接口调用节点 – 示例}
调用天气API获取实时温度，然后传入LLM节点生成穿衣建议。
\end{frame}

% 条件判断节点
\begin{frame}{IF节点 – 条件分支}
\begin{itemize}
\item 基于表达式（如 {{$json.value > 10}}）
\item 支持多分支（if-else if-else）
\item 可使用布尔运算符（and, or, not）
\end{itemize}
\end{frame}
\begin{frame}{IF节点 – 典型场景}
判断LLM返回的意图字段，分流到不同子工作流。
\end{frame}

% Merge节点
\begin{frame}{Merge节点 – 数据合并}
\begin{itemize}
\item 支持多种合并方式：
\begin{itemize}
\item Append：拼接数组
\item Wait：等待多个输入
\item Pass-through：传递第一个输入
\end{itemize}
\item 用于并行分支结果汇总
\end{itemize}
\end{frame}
\begin{frame}{Merge节点 – 示例}
同时调用两个API（用户信息、订单信息），Merge后生成完整报告。
\end{frame}

% Edit节点
\begin{frame}{Edit节点 – 字段操作}
\begin{itemize}
\item 添加/删除/重命名字段
\item 使用表达式计算新值
\item 支持JSONPath提取嵌套数据
\end{itemize}
\end{frame}
\begin{frame}{Edit节点 – 数据清洗}
去除无用字段，统一格式，为后续节点准备干净数据。
\end{frame}

% 循环节点
\begin{frame}{循环节点 – 类型}
\begin{itemize}
\item Split In Batches：将数组分批处理
\item Loop Over Items：对每个元素执行子工作流
\item 可设置并发数、超时
\end{itemize}
\end{frame}
\begin{frame}{循环节点 – 应用场景}
批量处理1000条用户记录，每条调用LLM生成个性化推荐。
\end{frame}
\begin{frame}{循环节点 – 终止条件}
可结合IF节点，满足条件时提前退出循环。
\end{frame}

% n8n高级特性
\begin{frame}{n8n + 扣子对比}
\begin{tabular}{p{0.4\textwidth}p{0.5\textwidth}}
\toprule
扣子 & n8n \\
\midrule
专为AI智能体设计 & 通用自动化平台 \\
闭源、托管 & 开源、自托管 \\
插件市场丰富 & 节点更多（400+） \\
知识库原生支持 & 需集成向量数据库 \\
适合快速构建对话Agent & 适合复杂系统集成 \\
\bottomrule
\end{tabular}
\end{frame}

\section{综合案例与总结}

\begin{frame}{案例：智能客服系统（扣子）}
\begin{enumerate}
\item 知识库：上传产品FAQ、操作手册
\item 插件：订单查询API、物流追踪
\item 工作流：用户意图分类 → 检索知识库 → 调用API → 生成回复
\item 发布到微信公众号
\end{enumerate}
\end{frame}

\begin{frame}{案例：自动化报表生成（n8n）}
\begin{enumerate}
\item 定时触发（每周一9点）
\item 从数据库查询销售数据
\item 数据分析节点计算周环比
\item LLM节点生成中文报告
\item 发送邮件给团队
\end{enumerate}
\end{frame}

\begin{frame}{平台选择建议}
\begin{itemize}
\item 快速构建对话Agent → 扣子、Dify
\item 复杂系统集成、数据处理 → n8n
\item AI原生开发 → Cursor、trae
\item 知识库问答 → Dokie、扣子
\item 完全自主可控 → n8n自托管 + 本地模型
\end{itemize}
\end{frame}

\begin{frame}{未来趋势}
\begin{itemize}
\item 智能体平台将融合低代码 + AI
\item 多智能体协作（A2A协议）
\item 边缘智能体（运行在手机、IoT设备）
\item 标准化工具接口（MCP）
\end{itemize}
\end{frame}

\begin{frame}{总结}
\begin{itemize}
\item 掌握扣子和n8n可覆盖绝大多数智能体开发需求
\item 低代码平台大幅降低开发门槛
\item 实际项目需结合具体场景选型
\end{itemize}
\end{frame}

\begin{frame}{Q\&A}
欢迎提问！
\vspace{1cm}
\textbf{资源链接：}
\begin{itemize}
\item 扣子官网：https://www.coze.cn
\item n8n文档：https://docs.n8n.io
\item Dify：https://dify.ai
\end{itemize}
\end{frame}
